\subsection{Controller 6: Adaptive Robust Gain (Asymptotic)}
\label{sec: controller6}

Controller 5-1 requires the robust gain to satisfy
$\rho_i \ge \|Y_i\tilde{\Theta}_i\| + \eta_i$, which presupposes knowledge of the uncertain
term $Y_i\tilde{\Theta}_i$. Controller 6 removes this requirement by letting the gain be
adapted online, at the price of losing the finite-time sliding property: $s_i$ now converges
to zero asymptotically (rather than in finite time), so the fencing objectives become
practical rather than asymptotic \cite{slotineLiAppliedNonlinearControl1991}.

\subsubsection{Integral APF and Sliding Variable}

As in Controller 5-1, define
\begin{equation}\label{eq6:zeta}
	\zeta_i = \dot{q}_0 - k_\alpha q_{i0}
	+ \int_0^t \bigl(-k_\alpha q_{i0} + \phi_i(\tau)\bigr)\,d\tau,
\end{equation}
so that
\begin{equation}\label{eq6:zeta_dot}
	\dot{\zeta}_i = \ddot{q}_0 - k_\alpha \dot{q}_{i0} - k_\alpha q_{i0} + \phi_i
\end{equation}
contains $\phi_i$ but not $\dot{\phi}_i$. The sliding variable is
\begin{equation}\label{eq6:sliding_var}
	s_i = \dot{q}_i - \zeta_i,
\end{equation}
with
\begin{equation}\label{eq6:s_dot_def}
	\dot{s}_i = \ddot{q}_{i0} + k_\alpha \dot{q}_{i0} + k_\alpha q_{i0} - \phi_i.
\end{equation}

\subsubsection{Control Law and Adaptation Laws}

The control input is
\begin{equation}\label{eq6:control_input}
	\tau_i = -k s_i - \hat{\rho}_i\,\mathrm{sgn}(s_i)
	+ Y_i(q_i,\dot{q}_i,\dot{\zeta}_i,\zeta_i) \hat{\Theta}_i,\qquad k>0,
\end{equation}
where $\hat{\rho}_i$ is an online estimate of the true bound
$\rho_i^* = \sup_t\|Y_i(t)\tilde{\Theta}_i(t)\|$. The adaptation laws are
\begin{align}
	\dot{\hat{\Theta}}_i &= -\Lambda^{-1} Y_i^T(q_i,\dot{q}_i,\dot{\zeta}_i,\zeta_i) s_i,
	\label{eq6:adaptation}\\
	\dot{\hat{\rho}}_i &= \gamma\,\|s_i\|_1,\qquad \hat{\rho}_i(0) > 0,
	\label{eq6:rho_adapt}
\end{align}
with $\Lambda > 0$ and $\gamma > 0$. The regression equation is
\begin{equation}\label{eq6:regression}
	M_i(q_i)\dot{\zeta}_i + C_i(q_i,\dot{q}_i)\zeta_i + g_i(q_i)
	= Y_i(q_i,\dot{q}_i,\dot{\zeta}_i,\zeta_i)\Theta_i.
\end{equation}

\subsubsection{Closed-Loop Dynamics and Lyapunov Analysis}

Substituting \eqref{eq6:control_input} into \eqref{eq hetero agent} and using
\eqref{eq6:regression} gives
\begin{equation}\label{eq6:s_dynamics}
	M_i(q_i)\dot{s}_i + C_i(q_i,\dot{q}_i)s_i
	= -k s_i - \hat{\rho}_i\,\mathrm{sgn}(s_i)
	+ Y_i(q_i,\dot{q}_i,\dot{\zeta}_i,\zeta_i)\tilde{\Theta}_i.
\end{equation}
Let $\tilde{\rho}_i = \hat{\rho}_i - \rho_i^*$ and consider
\begin{equation}\label{eq6:V}
	V_{1i} = \frac12 s_i^T M_i s_i
	+ \frac12 \tilde{\Theta}_i^T \Lambda \tilde{\Theta}_i
	+ \frac{1}{2\gamma}\tilde{\rho}_i^2.
\end{equation}
Differentiating and using \eqref{eq6:s_dynamics}, \eqref{eq6:adaptation},
\eqref{eq6:rho_adapt} and the skew-symmetry property A2,
\begin{align}
	\dot{V}_{1i}
	&= -k\|s_i\|^2 - \hat{\rho}_i\|s_i\|_1
		+ s_i^T Y_i\tilde{\Theta}_i
		- \tilde{\Theta}_i^T Y_i^T s_i
		+ \tilde{\rho}_i\|s_i\|_1 \notag\\
	&= -k\|s_i\|^2 - \rho_i^*\|s_i\|_1 \le 0. \label{eq6:V_dot}
\end{align}
Hence $V_{1i}$ is bounded, so $s_i \in \mathcal{L}_2 \cap \mathcal{L}_\infty$,
$\tilde{\Theta}_i$ bounded, and $\hat{\rho}_i$ bounded. Since $\dot{s}_i$ is uniformly
bounded (once all states are shown bounded), Barbalat's lemma gives $s_i \to 0$. No PE
condition is required.

\subsubsection{Theoretical Guarantee}

\begin{theorem}
	Let Assumptions A1--A3 hold. The HELS \eqref{eq hetero agent} governed by
	\eqref{eq6:control_input}, \eqref{eq6:adaptation} and \eqref{eq6:rho_adapt}
	achieves practical fencing: $s_i \to 0$ without the PE condition, and the fencing
	objectives P1--P3 hold in the practical sense, with ultimate bounds proportional to
	$\sup_t\|\dot{s}_i(t)\|$. No a~priori bound on $\|Y_i\tilde{\Theta}_i\|$ is needed.
\end{theorem}

\noindent\textbf{Proof}
The identity \eqref{eq6:V_dot} gives $s_i \to 0$ via Barbalat's lemma, with no PE
condition, as argued above.

\noindent\textit{Practical fencing.}
Because $s_i \to 0$ only asymptotically (not in finite time), $\dot{s}_i$ does not vanish
identically. From \eqref{eq6:s_dot_def},
$\ddot{q}_{i0} = -k_\alpha\dot{q}_{i0} - k_\alpha q_{i0} + \phi_i + \dot{s}_i$, and the
fencing Lyapunov function
\begin{equation*}
	V_2 = \frac12\sum_{i,j}\int_{\|q_{ij}\|}^{\mu}\alpha(s)\,ds
		+ \frac{k_\alpha}{2}\sum_i\|q_{i0}\|^2
		+ \frac12\sum_i\|\dot{q}_{i0}\|^2
\end{equation*}
satisfies
\begin{equation*}
	\dot{V}_2 = -k_\alpha\sum_i\|\dot{q}_{i0}\|^2 + \sum_i\dot{q}_{i0}^T\dot{s}_i
	\le -\frac{k_\alpha}{2}\sum_i\|\dot{q}_{i0}\|^2
		+ \frac{1}{2k_\alpha}\sum_i\|\dot{s}_i\|^2.
\end{equation*}
Since $s_i \to 0$ and $\dot{s}_i$ is bounded, the residual $\sum_i\|\dot{s}_i\|^2$ is
ultimately small. By the same Barbalat/ISS argument as in Controller 1, the centroid
error $\tilde{q}$ and the velocity errors $\|\dot{q}_{i0}\|$ are ultimately bounded by
quantities proportional to $\sup_t\|\dot{s}_i(t)\|$, establishing P1 and P3 in the
practical sense. Collision avoidance (P2) follows from the boundedness of $V_2$ on finite
intervals, exactly as in Controller 1.
\hfill$\square$

\begin{remark}
	The gain $\hat{\rho}_i$ grows monotonically according to \eqref{eq6:rho_adapt} and
	remains bounded because $s_i \in \mathcal{L}_1$ (from $\dot{V}_{1i} \le -k\|s_i\|^2$).
	No a~priori specification of $\rho_i$ is required. The price is that $s_i$ converges
	only asymptotically, so the finite-time sliding property of Controller 5-1 is lost
	and the fencing objectives become practical rather than asymptotic.
\end{remark}